\section{Estimating the energy on a fixed fiber}\label{sec:barrier}

On a fiber at fixed $(T,\tau)$, the only positive contribution to
the energy balance comes from the wave current crossing the spatial
wall $r=0$. The Gaussian is small where the current reaches the wall;
comparison with a spatial average of the wave bounds the other factor,
the boundary value of $u/r$. This gives an upper bound for the flux in
terms of the wave norm. To bound it by the energy itself, we must
recover that same norm from the interior. There we invert the Gaussian
and compare $u$ with a quadratic on each ball of radius $h$.

\begin{proposition}\label{prop:fixed-barrier}
Fix an annulus $c_1\le h|\xi|\le c_2$, with $0<c_1<c_2$. There are
constants $\gamma>0$, $K\ge1$, and $c,C>0$, depending only on $n$ and
the annulus, with the following property. Let $a=Kh$, $0<h<1/(4K)$,
and let $f$ be a trigonometric polynomial on the fixed torus, with
$\|f\|_2=1$ and frequencies in this annulus and in $|v|\le\gamma$ in
some orthonormal frame. Write $g_T=e^{-iT|D|}f$ and let
$\mathcal F(T,\tau)$ be its energy from Section~\ref{sec:energy}.
Then
\begin{align}
 \mathcal F(0,\tau)&\le C a^2\|f\|_{\pc}^{\pc}
       && (|\tau|=4a),\label{eq:barrier-input}\\
 \mathcal F(T,\tau)&\ge c a^2\|g_T\|_{\pc}^{\pc}
       && (|\tau|=2a),\label{eq:barrier-output}\\
 \Delta_z\mathcal F(T,\tau)&\le V(s)\mathcal F(T,\tau),\qquad
   V(s)=C s^{-2}e^{-c s/h}
       && (2a\le s=|\tau|\le3).\label{eq:barrier-potential}
\end{align}
The last two inequalities hold for all $T\in\mathbb R$.
\end{proposition}

\begin{proof}
Write $d=n-1$, $q=n/d=\pc/2$, and $\sigma=s/h$. Throughout the proof
$2a\le s\le3$. On the fiber
\[
 r=|W|^2/w_0^2+100|Q|^2-1<0,\qquad w_0=1/10,\quad Q=\Phi/a,
\]
we use two regions. On $r=0$ the cutoff in the current vanishes unless
$|Q|\le1/20$, which forces $|W|\ge (\sqrt{3}/2)w_0$. Near the center,
where $|W|\le\gamma$ and $|Q|\le1/400$, the Gaussian multiplier can
be inverted. Figure~\ref{fig:fiber} shows the separation between these
regions. It lets us compare two exponentials with different constants
in their exponents. For now take $0<\gamma<w_0/8$. We estimate uniformly in $s$,
allowing the prefactor in the interior lower bound to depend on $\gamma$.

\begin{figure}[!htbp]
\centering
\definecolor{fiberteal}{RGB}{27,88,96}
\definecolor{fiberclay}{RGB}{153,73,39}
\begin{tikzpicture}[
    font=\small,>=Latex,
    axis/.style={->,draw=black!72,line width=.5pt},
    guide/.style={draw=black!53,densely dotted,line width=.5pt},
    tick/.style={draw=black!70,line width=.5pt},
    inverse/.style={draw=fiberteal,line width=.8pt,dash pattern=on 3pt off 2pt},
    center patch/.style={draw=fiberteal,line width=.8pt},
    heading/.style={font=\small\bfseries,anchor=west,inner sep=0pt},
    every node/.style={inner sep=2pt}]
  % The horizontal unit is w_0, the vertical one 1/10.
  % Only the ratio gamma/w_0 is schematic.
  \begin{scope}[x=4.15cm,y=3.0cm]
    \node[heading] at (-.11,1.48) {(a) The fiber wall};
    \fill[black!3] (0,0) -- (1,0) arc[start angle=0,end angle=90,radius=1] -- cycle;
    \fill[fiberteal!7] (0,0) rectangle (.22,.1);
    \fill[fiberteal!28] (0,0) rectangle (.11,.025);
    \draw[inverse] (0,.1) -- (.22,.1) -- (.22,0);
    \draw[center patch] (0,.025) -- (.11,.025) -- (.11,0);
    \draw[guide] (0,.5) -- ({sqrt(3)/2},.5);
    \draw[black!78,line width=.65pt] (1,0) arc[start angle=0,end angle=90,radius=1];
    \draw[fiberclay,line width=2.05pt,line cap=round]
      (1,0) arc[start angle=0,end angle=30,radius=1];
    \draw[axis] (-.022,0) -- (1.16,0) node[right] {$|W|$};
    \draw[axis] (0,-.034) -- (0,1.15) node[above right] {$|Q|=|\Phi|/a$};
    \node[below left] at (0,0) {$0$};
    \draw[tick] (1,0) -- ++(0,-.025) node[below] {$w_0$};
    \draw[tick] (0,.5) -- ++(-.022,0) node[left] {$1/20$};
    \draw[tick] (0,1) -- ++(-.022,0) node[left] {$1/10$};
    \node at (.37,.71) {$r<0$};
    \node[anchor=west] at (.74,.76) {$r=0$};
    \node[fiberclay,anchor=west] (walltext) at (.37,.24) {flux wall};
    \draw[fiberclay,line width=.55pt] (walltext.east) -- (.80,.24) -- (.963,.257);
    \draw[fiberteal,line width=.5pt] (.18,.105) -- (.115,.22);
    \node[anchor=south,text=fiberteal,inner sep=1pt] at (.105,.225) {(b)};
  \end{scope}
  % Magnification of the lower left corner; the units are now 2 gamma
  % and 1/100. The axes in either panel are moduli, not wall pieces.
  \begin{scope}[xshift=6.85cm,x=3.65cm,y=2.8cm]
    \node[heading] at (-.14,1.586) {(b) The recovery patch};
    \fill[fiberteal!7] (0,0) rectangle (1,1);
    \fill[fiberteal!28] (0,0) rectangle (.5,.25);
    \draw[inverse] (0,1) -- (1,1) -- (1,0);
    \draw[center patch] (0,.25) -- (.5,.25) -- (.5,0);
    \draw[axis] (-.025,0) -- (1.17,0) node[right] {$|W|$};
    \draw[axis] (0,-.036) -- (0,1.235) node[above right] {$|Q|$};
    \node[below left] at (0,0) {$0$};
    \draw[tick] (.5,0) -- ++(0,-.026) node[below] {$\gamma$};
    \draw[tick] (1,0) -- ++(0,-.026) node[below] {$2\gamma$};
    \draw[tick] (0,.25) -- ++(-.025,0) node[left] {$1/400$};
    \draw[tick] (0,1) -- ++(-.025,0) node[left] {$1/100$};
    \node[align=center] at (.50,.66) {Gaussian inversion\\[-1pt]$g(Q)=1$};
    \node at (.25,.125) {centers};
  \end{scope}
\end{tikzpicture}
\caption{The fixed fiber in modulus coordinates. The wave flux through
$r=0$ is supported on the heavy arc, where $|Q|\le1/20$ and hence
$|W|\ge\sqrt{3}\,w_0/2$. This separation gives Gaussian decay
$e^{-cs/h}$ even for $|v|\le w_0/3$. In the enlargement (b), once
$K\gamma$ is large, the radius-$h$ balls about the marked centers
project into the dashed patch. There $|W-v|\le4\gamma$ on the enlarged cap
$|v|\le2\gamma$, so Gaussian inversion costs at most $e^{C\gamma^2s/h}$.
The aperture $\gamma$ is fixed independently of $h$; its size relative
to $w_0$ is schematic.}
\label{fig:fiber}
\end{figure}


\paragraph{A spatial average for the upper bound.}
For a dimensional constant $D\ge1$, to be fixed below, set
\[
 P_*(X)=\big((J_{Ds}*|g_T|^2)(X)\big)^{1/d},\qquad
 M=2n+4.
\]
Here and below convolution is in the real spatial variables $X$.
Recall that $J_\ell$ is the periodization of
$C_n\ell^{-n}(1+|X/\ell|^2)^{-M/2}$, of integral one. Differentiating this
positive kernel term by term and then taking the $d$th root gives
\begin{equation}\label{eq:barrier-average}
 |\nabla_X^jP_*|\le C(Ds)^{-j}P_*\quad(j=1,2),\qquad
 \int P_*^{d+1}\,dX\le\|g_T\|_{\pc}^{\pc}.
\end{equation}
The integral bound is Young's inequality with exponent $q$. The positive
kernel gives $P_*>0$ since $\|g_T\|_2=1$.

We next bound the wave current $A_{\mathrm w}$ and its regularized
density $\mu$. By Proposition~\ref{prop:current}, each scalar part of
the wave polarization is a multiplier on $g_T$ with symbol a bounded
smooth factor times
\[
 (\kappa/y)^{d/2}\exp(-\Psi),\qquad
 \Psi=\tfrac12\kappa y|W-v|^2,\qquad y=|Y|\asymp s.
\]
Multiply the symbol by a fixed smooth annular cutoff equal to one on the
data and supported in $|v|<w_0/3$. In the variable $\eta=h\xi$ its
derivatives through order $M$ are bounded by
$C(hs)^{-d/2}\sigma^m$ for some fixed $m$. Indeed $h\kappa$ and $v$
have uniformly bounded derivatives on this cutoff, $h\kappa$ is bounded
away from zero there, and each derivative of $\Psi$ costs at most
$C\sigma$. On $r=0$ where $g(Q)$ is nonzero, $|W-v|$ has a fixed
positive lower bound, so the same estimate gains $e^{-c_0\sigma}$ for
some $c_0>0$ independent of $\gamma$.

If $B$ bounds these derivatives of a compactly supported symbol, applying
$(1-\Delta_\eta)^{n+2}$ in its inverse Fourier transform bounds its
absolute kernel by $CB h^{-n}(1+|X/h|^2)^{-M/2}$. Periodization gives
the majorant $CBJ_h$ on the torus, also for a rotated lattice. Thus
the squared modulus of the multiplier output is bounded by
$C B^2 J_h*|g_T|^2$. There is one more convolution with $J_h$ in
$A_{\mathrm w}$. We have $J_h*J_h\le C J_h$: in the Euclidean
convolution, either $|Y|$ or $|X-Y|$ is at least $|X|/2$, and the
inequality persists after periodization. Moreover
$J_h\le(Ds/h)^nJ_{Ds}$. Using $(hs)^{-d}=s^{-2d}\sigma^d$ and
increasing $m$ when necessary, we obtain
\begin{align}
 \|A_{\mathrm w}\|&\le Cs^{-2d}\sigma^m P_*^d
              &&(r\le0),\nonumber\\
 \|A_{\mathrm w}\|&\le Cs^{-2d}\sigma^m e^{-c_0\sigma}P_*^d
              &&(r=0).\label{eq:barrier-current}
\end{align}
Here $\|A_{\mathrm w}\|$ is the matrix operator norm. Only the first
bound is needed for $\mu$. Its extra term is
$h^{50n}\det H_{\mathrm{fib}}(r)\le Ch^{50n}a^{-2}s^{-2d}$ by
Lemma~\ref{lem:fiber-geometry}. The nearest lattice term in the positive
kernel gives $\inf_X J_{Ds}(X)\ge c_D h^{M-n}$, since the torus has fixed
diameter and $h\le s\le3$. Consequently
$P_*^d\ge c_D h^{M-n}\|g_T\|_2^2=c_D h^{M-n}$, which absorbs
$h^{50n}a^{-2}$. In particular
\begin{equation}\label{eq:barrier-density}
 \mu\le C s^{-2d}\sigma^m P_*^d\qquad(r\le0).
\end{equation}

\paragraph{A comparison function that vanishes on the wall.}
Put $R=H_{\mathrm{fib}}(r)$. Two features of the defining function
from Lemma~\ref{lem:fiber-geometry} will be used together:
$R\ge c s^{-2}I$ and
$G_{\mathrm{fib}}(r)G_{\mathrm{fib}}(r)^*\le C R$. It follows that
\begin{equation}\label{eq:barrier-product}
 H_{\mathrm{fib}}(rP_*)\ge (1-C/D)P_*R.
\end{equation}
In fact, the product rule of Lemma~\ref{lem:hessian-calculus} gives
$P_*R$, $rH_{\mathrm{fib}}(P_*)$, and two terms pairing the gradients
of $r$ and $P_*$. For any quaternionic vector $e$ the two gradient
pairings with $e$ have moduli at most $C(e^*Re)^{1/2}$ and
$C D^{-1}P_*(e^*Re)^{1/2}$, respectively. This estimates the mixed
terms. The term involving $H_{\mathrm{fib}}(P_*)$ has relative cost
$C D^{-2}$ by \eqref{eq:barrier-average} and $|r|\le1$.

Since $r<0$, an upper bound on $u/r$ calls for a comparison function
below $u$, with determinant at least $\lambda\mu$. Fix $D$ so the
right side of \eqref{eq:barrier-product} is at least
$P_*R/2$. The background Hessian $B_{\mathrm{fib}}=\lambda
H_{\mathrm{fib}}(|\Phi|^2)$ has rank one and its positive eigenvalue
is comparable to $\lambda=h^{-30n}$. Hence for
$v_*=C_1\sigma^{m_1}rP_*$,
\[
 \det\big(B_{\mathrm{fib}}+H_{\mathrm{fib}}(v_*)\big)
 \ge c\lambda (C_1\sigma^{m_1}P_*s^{-2})^d\ge\lambda\mu
\]
by \eqref{eq:barrier-density} with $m_1,C_1$ large enough. The
potentials $v_*$ and $u$ have the same zero boundary data; Dirichlet
comparison (Theorem~\ref{thm:dirichlet}) yields $v_*\le u\le0$.
Division by $r<0$ and passage to the boundary give
\begin{equation}\label{eq:barrier-quotient}
 0\le u/r\le C_1\sigma^{m_1}P_*\qquad(r\le0).
\end{equation}

By Lemma~\ref{lem:energy-comparison}, $\mathcal F\le\int(-u)\mu$ on the
fiber. Its imaginary volume element is comparable to
$a^2s^{2d}\,dW\,dQ$ by Lemma~\ref{lem:fiber-geometry}. Combining
\eqref{eq:barrier-average}, \eqref{eq:barrier-density} and
\eqref{eq:barrier-quotient} gives
\begin{equation}\label{eq:barrier-upper-all}
 \mathcal F(T,\tau)\le Ca^2\sigma^m\|g_T\|_{\pc}^{\pc}.
\end{equation}
For the flux we need the \emph{full} gradient of $r$. The tangency in
Proposition~\ref{prop:current} gives $A_{\mathrm w}G(|\Phi|^2)=0$;
the remaining part, $G(|W|^2/w_0^2)$, has size $O(s^{-1})$ since
$W=P/Y$ and $y\asymp s$. Thus $G(r)^*A_{\mathrm w}G(r)\le
C s^{-2}\|A_{\mathrm w}\|$. The Jacobian estimate also gives
$\int\delta(r)\,dP\,dZ\le C a^2s^{2d}$ by coarea on the fixed
ellipsoid in $(W,Q)$. Applying the wall bound in
\eqref{eq:barrier-current} and \eqref{eq:barrier-quotient} to the
flux integral in Proposition~\ref{prop:energy-flux}, we conclude that
\begin{equation}\label{eq:barrier-flux-all}
 \Delta_z\mathcal F(T,\tau)
  \le C (a/s)^2\sigma^m e^{-c_0\sigma}\|g_T\|_{\pc}^{\pc}.
\end{equation}

\paragraph{Recovering the wave in the interior.}
Consider a center $(X,P,Z)$ with $|W|\le\gamma$ and $|Q|\le1/400$.
The ball of radius $h$ in the $3n$ real spatial coordinates about this
center lies in $|W|\le2\gamma$, $|Q|\le1/100$, provided $K\gamma$
and $K$ are sufficiently large: $W=P/Y$ changes by at most $Ch/s$
and $\Phi$ by at most $Ch$ there. The cutoff in $A_{\mathrm w}$ is
one throughout the ball.

The first scalar multiplier $\mathsf S_0$ of the time component in
\eqref{eq:current-time-multipliers} has symbol
$b_0(\kappa/y)^{d/2}e^{-\Psi}$ with
$b_0=(1+\operatorname{Re}(W)\cdot v/2)/\sqrt{2}$ bounded away from
zero. The other scalar component is orthogonal, so
$|\mathsf S_0g_T|^2\le A^0_{TT}$. Invert $\mathsf S_0$ on the frequencies of $g_T$
by multiplying its reciprocal symbol by a smooth annular cutoff equal
to one on the data and supported in $|v|\le2\gamma$. On its support
$\Psi\le C\gamma^2\sigma$. The derivatives of the reciprocal symbol in
$\eta=h\xi$ through order $M$ are bounded by
\[
 B_{\rm inv}=C_\gamma(hs)^{d/2}\sigma^m e^{C\gamma^2\sigma}.
\]
Here the derivatives of the cutoff enter only the prefactor $C_\gamma$;
the constant in the exponential is independent of $\gamma$.
The kernel estimate above yields
$|g_T(X)|^2\le C B_{\rm inv}^2 (J_h*A^0_{TT})(X)$ for each choice
of the imaginary coordinates in the ball. Since $J_h(X'-Y)$ and
$J_h(X-Y)$ are comparable for $|X'-X|\le h$, the density at
\emph{every point} in the ball is at least
\begin{equation}\label{eq:barrier-ball-density}
 m_X=c_\gamma(hs)^{-d}\sigma^{-m}e^{-C_2\gamma^2\sigma}|g_T(X)|^2.
\end{equation}

For the interior lower bound on $-u$ we need a comparison function
\emph{above} $u$. Let $l$ be Euclidean distance from the center. The quadratic
$v_b=\epsilon m_X^{1/d}(l^2-h^2)$ vanishes on the boundary of the ball
and has fiber Hessian $6\epsilon m_X^{1/d}I$ (three real slots per
spatial coordinate). For $m_X>0$, its determinant after adding
$B_{\mathrm{fib}}$ is at most
\[
 C\big((\epsilon m_X^{1/d})^n+\lambda(\epsilon m_X^{1/d})^d\big)
 \le \lambda m_X
\]
with a fixed small $\epsilon$: the density ceiling
\eqref{eq:energy-density-ceiling} and \eqref{eq:barrier-ball-density}
give $m_X^{1/d}\le\mu^{1/d}\le\lambda$. Now $u\le0=v_b$ on the ball
boundary and the determinant of $B_{\mathrm{fib}}+H_{\mathrm{fib}}(u)$
is $\lambda\mu\ge\lambda m_X$. The local comparison principle from
the proof of Theorem~\ref{thm:dirichlet} gives $u\le v_b$ in the ball.
At the center the depth is therefore at least $\epsilon h^2m_X^{1/d}$.
Multiplication by $\mu\ge m_X$ gives $(-u)\mu\ge\epsilon h^2m_X^{1+1/d}$,
also when $m_X=0$. Here $1+1/d=q$ and $m_X$ contains $|g_T(X)|^2$,
giving the required power $2q=\pc$ of the wave.

The imaginary volume of the center region is comparable to
$a^2s^{2d}\gamma^{2d}$. Integrate the center inequality there and in
$X$, use the lower bound for $\mathcal F$ in
Lemma~\ref{lem:energy-comparison}, and recall that
$h^2s^{2d}(hs)^{-dq}=\sigma^{n-2}$. Enlarging the fixed power $m$
and absorbing $\gamma^{2d}$ into the prefactor gives
\begin{equation}\label{eq:barrier-lower-all}
 \mathcal F(T,\tau)\ge
 c_\gamma a^2\sigma^{-m}e^{-C_3\gamma^2\sigma}\|g_T\|_{\pc}^{\pc}.
\end{equation}
Choose $\gamma$ so small that $C_3\gamma^2<c_0/2$, then fix $K$ to
meet the ball inclusions and to ensure that $h<1/(4K)$ lies in the
range of \eqref{eq:energy-density-ceiling}. Dividing the flux bound
\eqref{eq:barrier-flux-all} by the interior lower bound
\eqref{eq:barrier-lower-all} gives
\[
 \frac{\Delta_z\mathcal F}{\mathcal F}
   \le C_\gamma s^{-2}\sigma^{2m}e^{-(c_0-C_3\gamma^2)\sigma}.
\]
The exponential absorbs the fixed power of $\sigma$, proving
\eqref{eq:barrier-potential}. On the two circles in the $\tau$ plane
the scaled radii $4a/h=4K$ and $2a/h=2K$ are fixed. The energy
upper bound \eqref{eq:barrier-upper-all} at the first circle gives
\eqref{eq:barrier-input} for the initial data; the lower bound
\eqref{eq:barrier-lower-all} at the second gives \eqref{eq:barrier-output}
for the wave at time $T$.
\end{proof}
